\documentclass[10pt,a4paper]{amsart}
\usepackage[margin=19mm]{geometry}
\usepackage{amsmath,amssymb,mathtools}
\usepackage[scr=rsfs,cal=euler]{mathalfa}
\usepackage{xfrac}
\usepackage[unicode,hidelinks,bookmarks=false]{hyperref}
\newcommand{\ie}{i.e.\ }
\newcommand{\cf}{cf.\ }
\newcommand{\resp}{respectively\ }
\newcommand{\Subsection}{\subsection}
\providecommand{\nobreakdash}{\nobreak\hbox{-}}
\setlength{\emergencystretch}{2em}
\multlinegap0pt
\usepackage{bm}
\usepackage{bbm}
\usepackage{dsfont}
\usepackage{xspace}
\usepackage{cancel}
\usepackage{mathtools}
\usepackage{amsthm}
\makeatletter
\@ifundefined{theorem}{\newtheorem{theorem}{Theorem}}{}
\makeatother
\title[Two dimensional anisotropic mean curvature flow with contact]{Two dimensional anisotropic mean curvature flow with contact angle condition}
\author{Can Cui, Nung Kwan Yip}
\date{Source: arXiv:2510.22136v1 (CC BY 4.0)}
\begin{document}
\maketitle

\noindent\emph{Source and licence.} Two dimensional anisotropic mean curvature flow with contact angle condition. Version arXiv:2510.22136v1 (CC BY 4.0), distributed under the Creative Commons Attribution 4.0 International licence, as recorded in the arXiv licence metadata for this exact version and shown on the abstract page https://arxiv.org/abs/2510.22136v1. Full rights notice: licenses/arxiv-2510.22136-CC-BY-4.0.txt.\par\smallskip

\noindent\textit{Editorial scope.} Counted unit: the Theorem whose source label is \texttt{None} in the article (source0.tex lines 626--629, source label \texttt{None}), delivered together with the article's own complete original proof of it (source0.tex lines 630--672, one \texttt{proof} environment of 43 lines). No mathematics is written, repaired or re-derived: statement and proof are the article's own text line for line. Required context retained uncounted: eq:ell\_ang\_pro (lines 173--178). Every definition, display and label of those lines is retained unchanged. Omitted from this package and not redistributed: 1--172, 179--625, 673--997. Nothing inside a retained excerpt is omitted or altered: every retained body is delivered exactly as the article's lines are written, so no omission receipt of any kind accompanies this package. Citations: the retained excerpts carry no citation command at all, so no bibliography entry of the article is transcribed and no reference is redistributed. Reference adaptation: none was needed and none was made; every reference inside the delivered unit resolves inside it, so no unresolved reference or citation is produced. Printed slips kept literal: no printed word, symbol, hypothesis or number of the counted statement or of its proof was altered. Layout and class: the article's own class is \texttt{article} with packages including graphicx, amsmath, amssymb, amsthm, amsrefs, enumitem, xcolor, tikz, hyperref, lineno, geometry, authblk; the wrapper is the lane's pinned \texttt{amsart} editorial wrapper at 10pt on A4, which restates the theorem-like environments the retained text needs and adds the article's own macro definitions that the retained text needs (none), with extra wrapper packages (bm, bbm, dsfont, xspace, cancel, mathtools). The delivered numbering is the unit's own. Assets: no retained span contains a figure environment, a graphics-inclusion command, a PDF-page inclusion, a table or a TikZ picture; no image file is copied into this package, the package declares no asset dependency and no image file is redistributed with it. Licence: version arXiv:2510.22136v1 of this article is distributed under the Creative Commons Attribution 4.0 International licence, as recorded in the arXiv licence metadata for this exact version and shown on the abstract page https://arxiv.org/abs/2510.22136v1. A full rights notice is delivered at licenses/arxiv-2510.22136-CC-BY-4.0.txt. Characters: every retained excerpt is pure ASCII, so the delivered engine, XeLaTeX with its default Latin Modern text font, reports no missing character.
\begin{equation} \label{eq:ell_ang_pro}
\left\{ \begin{array}{cllcl}
     \lambda & = & G (Dw,-1) \ D^2_{p_i p_j} F (Dw,-1) w_{x_i x_j}   \quad & \text{in} & \ \Omega,  \\
     D_N w & = & - \sqrt{1+|Dw|^2} \cos \theta \quad & \text{on} & \ \partial \Omega. 
\end{array} \right.
\end{equation}

\begin{theorem}
Under the assumptions in Theorem 3.1, there is a unique $\lambda$ such that a solution exists for \eqref{eq:ell_ang_pro}.
In this case, the solution is also unique.
\end{theorem}

\begin{proof}
Firstly, we want to show that there exists a positive constant $C$ independent of $\epsilon$ such that $|\epsilon w^{\epsilon}|\leq C$. 

For this purpose, consider the function $\psi=a_0 d$ defined in a neighborhood of $\partial\Omega$.
Here $d$ is the distance function to $\partial\Omega$ and $a_0$ is a constant such that 
\begin{equation*}
    a_0 < -\cos{\theta(\cdot)} \sqrt{1+a_0^2}.
\end{equation*}
(By assumption A2, such an $a_0$ exists.)
Then we extend $\psi$ smoothly to all of $\Omega$. Note that
$D_N\psi = a_0$ on $\partial\Omega$. Now suppose $\psi-w_{\epsilon}$ attains its minimum at $x_0$. We have following two cases: \\
\textbf{Case 1:} $x_0 \in \partial\Omega$. This implies $D_T \psi(x_0) = D_T w^{\epsilon}(x_0)$ and $D_N \psi(x_0) \geq D_N w^{\epsilon}(x_0)$. Hence at $x_0$, we have
\begin{equation*}
      -\cos{\theta} = \frac{D_N w^{\epsilon}}{\sqrt{1+|D w^{\epsilon}}|^2} \leq \frac{D_N \psi}{\sqrt{1+|D \psi|^2}} < -\cos{\theta},
\end{equation*}
which is a contradiction. (In the above, we have used the facts that 
$|Dw^\epsilon|^2 = |D_Tw^\epsilon|^2 + |D_Nw^\epsilon|^2
\leq
|D_T\psi|^2 + |D_N\psi|^2
=|D\psi|^2
$ and the function $f(x) = \frac{x}{\sqrt{c^2+x^2}}$ is increasing 
in $x\in\mathbb{R}$.)\\
\textbf{Case 2:} $x_0 \in \Omega$. This implies $D\psi(x_0) = D w^{\epsilon}(x_0)$ and $\big[ D^2 \psi(x_0) \big] \geq \big[ D^2 w^{\epsilon}(x_0) \big]$. Hence at $x_0$, for some constant $C=C(\theta, \Omega)$,
\begin{equation*}
    \epsilon w^{\epsilon} = a^{ij} (D w^{\epsilon}) w_{x_i x_j}^{\epsilon} \leq a^{ij} (D \psi) \psi_{x_i x_j} \leq C,
\end{equation*}
which suggests an upper bound for $\epsilon w^{\epsilon}$. Similarly, considering the maximum of $\psi-w_{\epsilon}$ suggests a lower bound. Hence $|\epsilon w^{\epsilon}|$ is bounded.

Once the above is achieved, we can then directly apply the proof of Theorem 3.1 by replacing $u_t$ with $\epsilon w^{\epsilon}$ to obtain an uniform estimate for $|Dw^{\epsilon}|$.
This implies $| D (\epsilon w^{\epsilon})| \rightarrow 0$ as $\epsilon \rightarrow 0$. Hence $\epsilon w^{\epsilon} \rightarrow \lambda$ for some constant $\lambda$.

For the uniqueness, assume $u^1$ and $u^2$ are two solutions of \eqref{eq:ell_ang_pro} with $\lambda_1$ and $\lambda_2$ respectively and $\lambda_1 \leq \lambda_2$. Let $\Tilde{u}=u^1 - u^2$. Then we have
\begin{equation*}
    \Tilde{a}^{ij} \Tilde{u}_{x_i x_j} + \Tilde{b}^{i} \Tilde{u}_{x_i} \leq 0,
\end{equation*}
where $\Tilde{a}^{ij} = a^{ij} (Du^1)$ and $\Tilde{b}^{i} = u^2_{x_k x_l} \int^1_0 a^{kl}_{p_i} (\eta Du^1 + (1-\eta) Du^2 ) d\eta $. The maximum principle suggests that 
$\max \Tilde{u}$ occurs on the boundary, $x_0 \in \partial \Omega$. Then Hopf lemma implies that at $x_0$, we have $D_N u^1 < D_N u^2$. Together with $D_T u^1 = D_T u^2$, this would then contradict the following,  
\begin{equation*}
    \frac{D_N u^1}{\sqrt{1+|D u^1|^2}} =  \frac{D_N u^2}{\sqrt{1+|D u^2|^2}}
    \,\,(=\cos\theta).
\end{equation*}
Hence $\Tilde{u}$ must be a constant and $\lambda$ is unique.
\end{proof} 

\end{document}
